Segundo \citeonline{SilvaVilela2021}, Machine Learning pode ser 
definida como um ramo da \gls{ia} responsável por capacitar computadores
a simular aprendizado com base em processamento de dados. 
Essa técnica está em ascensão atualmente, sendo amplamente aplicada em softwares que 
têm por objetivo a identificação de padrões, como é o caso do \gls{ocr},
detecção de phishing em e-mails, detecção de uso fraudulento de cartões de crédito, 
entre outras aplicações.

Tal subconjunto tem por característica a distinção de métodos de aprendizagem, 
tendo como os mais difundidos os métodos supervisionado, 
não supervisionado e semissupervisionado.

\subsubsubsection{Aprendizado Supervisionado}

Neste método de aprendizagem, o objetivo é ensinar ao modelo de \gls{ia}, 
através de objetos previamente rotulados, a identificar características de modo a 
classificar novos objetos \cite{SilvaVilela2021}. 
Entre os algoritmos mais importantes para a aplicação deste tipo de 
método de treinamento estão:

\begin{enumerate}
	\item k-Nearest Neighbours;
    \item Regressão Linear;
    \item Máquina de Vetores de Suporte (SVM);
    \item Árvore de Decisão e Floresta Aleatória;
    \item Redes Neurais;
\end{enumerate}

Observação: estes algoritmos não necessariamente são concorrentes; 
portanto, em um único modelo, pode haver implementação de um ou mais 
dos algoritmos citados \cite{SilvaVilela2021}.

\subsubsubsection{Deep Learning}
Deep Learning, ou aprendizado profundo, é uma técnica que se encontra 
em um subcampo do Machine Learning, cujo desenvolvimento foi pautado 
na simulação do processo de aprendizagem humana. 
Nesse caso, utiliza-se como base o algoritmo de redes neurais, 
construindo estruturas conectadas de forma a emular o comportamento do cérebro humano 
\cite{jesus2023inteligencia}.
